% Blank variant of src/architecture/decisions/templates/adr-observability-template.tex. Unapproved working document.
% --------------------------------------------------------------------------
% File: src/architecture/decisions/templates/adr-observability-template.tex
% Style: style v1.0 (tooling/latex/style.sty)
% Engine: pdfLaTeX through tooling/scripts/latex_build.py
% Purpose: ready-to-fill blank record; not an approved decision
% --------------------------------------------------------------------------
\documentclass[11pt,letterpaper]{article}
\usepackage{technical-design-spec}
\renewcommand{\DocTitle}{ADR: Observability, Telemetry and Operational Detection}
\renewcommand{\DocSubtitle}{Fillable working record and evidence trail}
\renewcommand{\DocVersion}{v1.0}
\renewcommand{\DocStatus}{Blank}
\renewcommand{\DocOwner}{Architecture Governance}
\renewcommand{\DocAuthor}{Architecture Documentation Library}
\renewcommand{\DocDate}{2026-10-28}
\renewcommand{\DocSummary}{An unfilled, domain-specific architecture decision record for observability, telemetry and operational detection. Inherits the central decision structure and adds specialized design evidence.}
% Printed forms deliberately contain no bracketed response placeholders.
\newcommand{\BlankResponse}{\rule[-0.4em]{0pt}{2em}\rule{0.65\linewidth}{0.4pt}}
\begin{document}
\makedoctocpage


\section{Record identity and decision authority}
\noindent\textbf{Status:} \rule{1.1in}{0.4pt}\quad
\textbf{ADR ID:} \rule{1.0in}{0.4pt}\quad
\textbf{Date:} \rule{1.0in}{0.4pt}

\begin{description}
\item[Decision title:] \BlankResponse
\item[Entity of interest:] \BlankResponse
\item[Scope and baseline:] \BlankResponse
\item[Owner and accountable approver:] \BlankResponse
\item[Consulted stakeholders:] \BlankResponse
\item[Related records:] \BlankResponse
\item[Review and expiry trigger:] \BlankResponse
\end{description}

\begin{notebox}[Unapproved working record]
Complete every required field. Acceptance requires recorded decision
authority and evidence; this Blank is never an approved decision.
\end{notebox}

\section{Context and architectural significance}
\begin{description}
\item[Problem and forces:] \BlankResponse
\item[Stakeholders and concerns:] \BlankResponse
\item[Drivers and constraints:] \BlankResponse
\item[Boundary and assumptions:] \BlankResponse
\item[Decision deadline:] \BlankResponse
\end{description}

\section{Options considered and evidence}
Include the status quo as an explicit alternative where practical. Evaluate
at least two feasible options, or explain why only one remains viable.
\begin{tabularx}{\linewidth}{@{}p{0.18\linewidth}X X@{}}
\toprule
\textbf{Option} & \textbf{Benefits and supporting evidence} & \textbf{Costs, risks and reason for disposition}\\
\midrule
A: \BlankResponse & \BlankResponse & \BlankResponse\\
B: \BlankResponse & \BlankResponse & \BlankResponse\\
C: \BlankResponse & \BlankResponse & \BlankResponse\\
\bottomrule
\end{tabularx}

\noindent\textbf{Evaluation method:} \BlankResponse

\section{Decision, rationale and non-goals}
\begin{description}
\item[We choose:] \BlankResponse
\item[Because:] \BlankResponse
\item[Non-goals:] \BlankResponse
\item[Invariants and exceptions:] \BlankResponse
\item[Dependencies:] \BlankResponse
\end{description}

\section{Domain-specific decision inputs}
Complete the following prompts with referenceable designs, measurements,
threat assessments or operational constraints; mark inapplicable prompts
explicitly rather than removing necessary decision context.
\begin{itemize}
\item Critical user journeys, golden signals, service-level indicators and SLO ownership.
\item Metrics/logs/traces topology, context propagation, structured events, sampling and instrumentation standards.
\item Alert thresholds, alert routes, error-budget burn policies, paging noise and incident escalation.
\item Sensitive telemetry handling, retention, regional residency, role-based access and audit requirements.
\item Telemetry volume, cardinality budgets, egress, vendor lock-in, data quality and dashboard ownership.
\end{itemize}
\section{Domain-specific alternatives and trade-offs}
Explain how each feasible option addresses the following criteria.
\begin{itemize}
\item Coverage, detection delay, diagnostic accuracy, cardinality and ongoing cost.
\item Centralized versus distributed collectors, vendor-native versus open instrumentation and resilience during outages.
\item Avoidance of sensitive data leaks, false positives and unowned alerts.
\end{itemize}
\noindent\textbf{Domain comparison:} \BlankResponse
\section{Domain-specific validation and review gates}
The proposed decision is ready for acceptance only after its applicable
tests have objective pass/fail thresholds, an owner and dated evidence.
\begin{itemize}
\item Synthetic failure produces correctly routed alerts within target detection latency.
\item Trace linkage across service boundaries; sampling/retention test and sensitive-field redaction.
\item Cost projections with cardinality stress, dashboard verification and incident drill.
\end{itemize}
\noindent\textbf{Related repository material:} Cloud service-quality indicators, monitoring materials and security logging controls.

\section{Consequences, risks and safeguards}
\begin{description}
\item[Positive consequences:] \BlankResponse
\item[Negative consequences:] \BlankResponse
\item[Security and privacy effects:] \BlankResponse
\item[Operational and financial effects:] \BlankResponse
\item[Unknowns and risk treatment:] \BlankResponse
\end{description}

\section{Implementation, verification and reversibility}
\begin{tabularx}{\linewidth}{@{}p{0.22\linewidth}X@{}}
\toprule
\textbf{Control} & \textbf{Complete with concrete, reviewable evidence}\\
\midrule
Implementation owner and plan & \BlankResponse\\
Acceptance tests & \BlankResponse\\
Security / reliability gates & \BlankResponse\\
Evidence and provenance & \BlankResponse\\
Rollback or exit & \BlankResponse\\
Post-deployment review & \BlankResponse\\
\bottomrule
\end{tabularx}

\section{Architecture-description traceability}
A decision record contributes rationale to an architecture description; it
does not, by itself, constitute a conforming ISO/IEC/IEEE 42010:2022
architecture description.
\begin{description}
\item[Description and configuration baseline:] \BlankResponse
\item[Stakeholders and concerns:] \BlankResponse
\item[Viewpoints, views and model kinds:] \BlankResponse
\item[Correspondence and consistency:] \BlankResponse
\item[Requirements and evidence:] \BlankResponse
\end{description}

\section{Review history and references}
\begin{tabularx}{\linewidth}{@{}p{0.16\linewidth}p{0.19\linewidth}p{0.19\linewidth}X@{}}
\toprule
\textbf{Date} & \textbf{Reviewer / role} & \textbf{Disposition} & \textbf{Evidence and requested changes}\\
\midrule
\emph{[date]} & \BlankResponse & \BlankResponse & \BlankResponse\\
\bottomrule
\end{tabularx}

\noindent\textbf{Sources:} \BlankResponse

\end{document}
